\begin{thebibliography}{99}

\bibitem{Alarcao2019} Alarc\~ao SM, Fonseca MJ (2019) Emotions recognition using EEG signals: a survey. IEEE Trans Affect Comput 10(3):374--393. \url{https://doi.org/10.1109/TAFFC.2017.2714671}

\bibitem{Arevalo2020} Arevalo J, Solorio T, Montes-y-G\'omez M, Gonz\'alez FA (2020) Gated multimodal networks. Neural Comput Appl 32(14):10209--10228. \url{https://doi.org/10.1007/s00521-019-04559-1}

\bibitem{Bandt2002} Bandt C, Pompe B (2002) Permutation entropy: a natural complexity measure for time series. Phys Rev Lett 88(17):174102. \url{https://doi.org/10.1103/PhysRevLett.88.174102}

\bibitem{Bradley2008} Bradley MM, Miccoli L, Escrig MA, Lang PJ (2008) The pupil as a measure of emotional arousal and autonomic activation. Psychophysiology 45(4):602--607. \url{https://doi.org/10.1111/j.1469-8986.2008.00654.x}

\bibitem{ChoKim1995} Cho SB, Kim JH (1995) Combining multiple neural networks by fuzzy integral for robust classification. IEEE Trans Syst Man Cybern 25(2):380--384. \url{https://doi.org/10.1109/21.364825}

\bibitem{Demsar2006} Dem\v{s}ar J (2006) Statistical comparisons of classifiers over multiple data sets. J Mach Learn Res 7:1--30

\bibitem{Duan2013} Duan RN, Zhu JY, Lu BL (2013) Differential entropy feature for EEG-based emotion classification. In: 6th Int IEEE/EMBS Conf Neural Engineering (NER), pp 81--84. \url{https://doi.org/10.1109/NER.2013.6695876}

\bibitem{FarhadiSedehi2024a} Farhadi Sedehi J, Jafarnia Dabanloo N, Maghooli K, Sheikhani A (2024) Recognition of emotional categories using mined fuzzy rules from electroencephalogram signals with gated eye track data approach. IEEE Access 12:118122--118140. \url{https://doi.org/10.1109/ACCESS.2024.3430934}

\bibitem{FarhadiSedehi2024b} Farhadi Sedehi J, Jafarnia Dabanloo N, Maghooli K, Sheikhani A (2024) Multimodal insights into Granger causality connectivity: integrating physiological signals and gated eye-tracking data for emotion recognition using convolutional neural network. Heliyon 10(16):e36411. \url{https://doi.org/10.1016/j.heliyon.2024.e36411}

\bibitem{Fu2023} Fu B, Gu C, Fu M, Xia Y, Liu Y (2023) A novel feature fusion network for multimodal emotion recognition from EEG and eye movement signals. Front Neurosci 17:1234162. \url{https://doi.org/10.3389/fnins.2023.1234162}

\bibitem{Hollenstein2018} Hollenstein N, Rotsztejn J, Troendle M, Pedroni A, Zhang C, Langer N (2018) ZuCo, a simultaneous EEG and eye-tracking resource for natural sentence reading. Sci Data 5:180291. \url{https://doi.org/10.1038/sdata.2018.291}

\bibitem{Hollenstein2020} Hollenstein N, Troendle M, Zhang C, Langer N (2020) ZuCo 2.0: a dataset of physiological recordings during natural reading and annotation. In: Proc 12th Language Resources and Evaluation Conf (LREC), pp 138--146

\bibitem{Holmqvist2012} Holmqvist K, Nystr\"om M, Mulvey F (2012) Eye tracker data quality: what it is and how to measure it. In: Proc Symp Eye Tracking Research and Applications (ETRA), pp 45--52. \url{https://doi.org/10.1145/2168556.2168563}

\bibitem{Jacobs1991} Jacobs RA, Jordan MI, Nowlan SJ, Hinton GE (1991) Adaptive mixtures of local experts. Neural Comput 3(1):79--87. \url{https://doi.org/10.1162/neco.1991.3.1.79}

\bibitem{Jang1993} Jang JSR (1993) ANFIS: adaptive-network-based fuzzy inference system. IEEE Trans Syst Man Cybern 23(3):665--685. \url{https://doi.org/10.1109/21.256541}

\bibitem{Jenke2014} Jenke R, Peer A, Buss M (2014) Feature extraction and selection for emotion recognition from EEG. IEEE Trans Affect Comput 5(3):327--339. \url{https://doi.org/10.1109/TAFFC.2014.2339834}

\bibitem{Kaya2024} Kaya \"U, Akay D, \c{S}eng\"ul Ayan S (2024) EEG-based emotion recognition in neuromarketing using fuzzy linguistic summarization. IEEE Trans Fuzzy Syst 32(8):4248--4259. \url{https://doi.org/10.1109/TFUZZ.2024.3392495}

\bibitem{Khan2025} Khan NN, Sweet T, Harvey CA, Knapp C, Krusienski DJ, Thompson DE (2025) The role of review process failures in affective state estimation: an empirical investigation of DEAP dataset. arXiv:2508.02417 (preprint)

\bibitem{Kittler1998} Kittler J, Hatef M, Duin RPW, Matas J (1998) On combining classifiers. IEEE Trans Pattern Anal Mach Intell 20(3):226--239. \url{https://doi.org/10.1109/34.667881}

\bibitem{Koelstra2012} Koelstra S, M\"uhl C, Soleymani M, Lee JS, Yazdani A, Ebrahimi T, Pun T, Nijholt A, Patras I (2012) DEAP: a database for emotion analysis using physiological signals. IEEE Trans Affect Comput 3(1):18--31. \url{https://doi.org/10.1109/T-AFFC.2011.15}

\bibitem{Kret2019} Kret ME, Sjak-Shie EE (2019) Preprocessing pupil size data: guidelines and code. Behav Res Methods 51(3):1336--1342. \url{https://doi.org/10.3758/s13428-018-1075-y}

\bibitem{Kukhilava2025} Kukhilava N, Tsmindashvili T, Kalandadze R, et al (2025) Evaluation in EEG emotion recognition: state-of-the-art review and unified framework. arXiv:2505.18175 (preprint)

\bibitem{Li2021} Li R, Johansen JS, Ahmed H, Ilyevsky TV, Wilbur RB, Bharadwaj HM, Siskind JM (2021) The perils and pitfalls of block design for EEG classification experiments. IEEE Trans Pattern Anal Mach Intell 43(1):316--333. \url{https://doi.org/10.1109/TPAMI.2020.2973153}

\bibitem{Lim2020} Lim JZ, Mountstephens J, Teo J (2020) Emotion recognition using eye-tracking: taxonomy, review and current challenges. Sensors 20(8):2384. \url{https://doi.org/10.3390/s20082384}

\bibitem{LinWang2002} Lin CF, Wang SD (2002) Fuzzy support vector machines. IEEE Trans Neural Netw 13(2):464--471. \url{https://doi.org/10.1109/72.991432}

\bibitem{Liu2022} Liu W, Qiu JL, Zheng WL, Lu BL (2022) Comparing recognition performance and robustness of multimodal deep learning models for multimodal emotion recognition. IEEE Trans Cogn Dev Syst 14(2):715--729. \url{https://doi.org/10.1109/TCDS.2021.3071170}

\bibitem{Lu2015} Lu Y, Zheng WL, Li B, Lu BL (2015) Combining eye movements and EEG to enhance emotion recognition. In: Proc 24th Int Joint Conf Artificial Intelligence (IJCAI), pp 1170--1175

\bibitem{Mallat1989} Mallat SG (1989) A theory for multiresolution signal decomposition: the wavelet representation. IEEE Trans Pattern Anal Mach Intell 11(7):674--693. \url{https://doi.org/10.1109/34.192463}

\bibitem{Mathot2018} Math\^ot S (2018) Pupillometry: psychology, physiology, and function. J Cogn 1(1):16. \url{https://doi.org/10.5334/joc.18}

\bibitem{Matiko2014} Matiko JW, Beeby SP, Tudor J (2014) Fuzzy logic based emotion classification. In: Proc IEEE Int Conf Acoustics, Speech and Signal Processing (ICASSP), pp 4389--4393. \url{https://doi.org/10.1109/ICASSP.2014.6854431}

\bibitem{McNemar1947} McNemar Q (1947) Note on the sampling error of the difference between correlated proportions or percentages. Psychometrika 12(2):153--157. \url{https://doi.org/10.1007/BF02295996}

\bibitem{Murugappan2010} Murugappan M, Ramachandran N, Sazali Y (2010) Classification of human emotion from EEG using discrete wavelet transform. J Biomed Sci Eng 3(4):390--396. \url{https://doi.org/10.4236/jbise.2010.34054}

\bibitem{Nandakumar2008} Nandakumar K, Chen Y, Dass SC, Jain AK (2008) Likelihood ratio-based biometric score fusion. IEEE Trans Pattern Anal Mach Intell 30(2):342--347. \url{https://doi.org/10.1109/TPAMI.2007.70796}

\bibitem{Partala2003} Partala T, Surakka V (2003) Pupil size variation as an indication of affective processing. Int J Hum-Comput Stud 59(1--2):185--198. \url{https://doi.org/10.1016/S1071-5819(03)00017-X}

\bibitem{Pedroni2019} Pedroni A, Bahreini A, Langer N (2019) Automagic: standardized preprocessing of big EEG data. NeuroImage 200:460--473. \url{https://doi.org/10.1016/j.neuroimage.2019.06.046}

\bibitem{Peng2005} Peng H, Long F, Ding C (2005) Feature selection based on mutual information criteria of max-dependency, max-relevance, and min-redundancy. IEEE Trans Pattern Anal Mach Intell 27(8):1226--1238. \url{https://doi.org/10.1109/TPAMI.2005.159}

\bibitem{Raiturkar2016} Raiturkar P, Kleinsmith A, Keil A, Banerjee A, Jain E (2016) Decoupling light reflex from pupillary dilation to measure emotional arousal in videos. In: Proc ACM Symp Applied Perception (SAP), pp 89--96. \url{https://doi.org/10.1145/2931002.2931009}

\bibitem{Richman2000} Richman JS, Moorman JR (2000) Physiological time-series analysis using approximate entropy and sample entropy. Am J Physiol Heart Circ Physiol 278(6):H2039--H2049. \url{https://doi.org/10.1152/ajpheart.2000.278.6.H2039}

\bibitem{Rohila2019} Rohila A, Sharma A (2019) Phase entropy: a new complexity measure for heart rate variability. Physiol Meas 40(10):105006. \url{https://doi.org/10.1088/1361-6579/ab499e}

\bibitem{Socher2013} Socher R, Perelygin A, Wu J, Chuang J, Manning CD, Ng A, Potts C (2013) Recursive deep models for semantic compositionality over a sentiment treebank. In: Proc Conf Empirical Methods in Natural Language Processing (EMNLP), pp 1631--1642. \url{https://doi.org/10.18653/v1/D13-1170}

\bibitem{Soleymani2012db} Soleymani M, Lichtenauer J, Pun T, Pantic M (2012) A multimodal database for affect recognition and implicit tagging. IEEE Trans Affect Comput 3(1):42--55. \url{https://doi.org/10.1109/T-AFFC.2011.25}

\bibitem{Soleymani2012mm} Soleymani M, Pantic M, Pun T (2012) Multimodal emotion recognition in response to videos. IEEE Trans Affect Comput 3(2):211--223. \url{https://doi.org/10.1109/T-AFFC.2011.37}

\bibitem{TakagiSugeno1985} Takagi T, Sugeno M (1985) Fuzzy identification of systems and its applications to modeling and control. IEEE Trans Syst Man Cybern SMC-15(1):116--132. \url{https://doi.org/10.1109/TSMC.1985.6313399}

\bibitem{WangMendel1992} Wang LX, Mendel JM (1992) Generating fuzzy rules by learning from examples. IEEE Trans Syst Man Cybern 22(6):1414--1427. \url{https://doi.org/10.1109/21.199466}

\bibitem{Winkler2011} Winkler I, Haufe S, Tangermann M (2011) Automatic classification of artifactual ICA-components for artifact removal in EEG signals. Behav Brain Funct 7:30. \url{https://doi.org/10.1186/1744-9081-7-30}

\bibitem{Zadeh1965} Zadeh LA (1965) Fuzzy sets. Inf Control 8(3):338--353. \url{https://doi.org/10.1016/S0019-9958(65)90241-X}

\bibitem{Zebrowski2000} \.Zebrowski JJ, Pop\l awska W, Baranowski R, Buchner T (2000) Symbolic dynamics and complexity in a physiological time series. Chaos Solitons Fractals 11(7):1061--1075. \url{https://doi.org/10.1016/S0960-0779(99)00004-1}

\bibitem{ZhengDong2014} Zheng WL, Dong BN, Lu BL (2014) Multimodal emotion recognition using EEG and eye tracking data. In: Proc 36th Annu Int Conf IEEE Eng Med Biol Soc (EMBC), pp 5040--5043. \url{https://doi.org/10.1109/EMBC.2014.6944757}

\bibitem{Zheng2019} Zheng WL, Liu W, Lu Y, Lu BL, Cichocki A (2019) EmotionMeter: a multimodal framework for recognizing human emotions. IEEE Trans Cybern 49(3):1110--1122. \url{https://doi.org/10.1109/TCYB.2018.2797176}

\bibitem{ZhengLu2015} Zheng WL, Lu BL (2015) Investigating critical frequency bands and channels for EEG-based emotion recognition with deep neural networks. IEEE Trans Auton Ment Dev 7(3):162--175. \url{https://doi.org/10.1109/TAMD.2015.2431497}

\end{thebibliography}
